\section{Autoencoders}
Autoencoders (AEs) are a class of neural networks trained in an unsupervised manner to learn efficient, compressed representations of data. 

\paragraph{The Core Idea}
Conceptually, autoencoders generalize the idea behind Principal Component Analysis (PCA): both aim to learn a lower-dimensional representation of the input data, known as the \textbf{latent space}.

While PCA is restricted to discovering strictly linear transformations, autoencoders leverage non-linear activation functions, allowing them to model much more complex, non-linear data manifolds. In fact, if an autoencoder uses only linear activations and is trained with a mean squared error loss, it learns to span the same principal subspace as PCA.

This lower-dimensional representation is learned by training the network to map the input data $x$ to a compressed space $z$ via an \textit{encoder}, and then reconstruct the original data as accurately as possible through a second network called the \textit{decoder}.

\begin{figure}[H]
    \centering
    \resizebox{0.6\textwidth}{!}{\usetikzlibrary{arrows.meta, calc}

\begin{tikzpicture}[
    >=Stealth,
    flow/.style={-{Stealth[length=6pt]}, line width=1pt},
    skip/.style={-{Stealth[length=6pt]}, line width=1pt, color=skipcolor},
    lbl/.style={font=\small, align=center},
    blocklabel/.style={font=\bfseries\small},
    latentblock/.style={rectangle, draw, thick, fill=dlgray!25, minimum width=1.4cm, minimum height=0.7cm, font=\small, align=center},
]
    % Geometry, defined explicitly (rather than relying on the trapezium
    % shape's anchors) so the loss arrow can be routed with exact clearance.
    % Both blocks share the wide/narrow edge widths, so the silhouette forms
    % a symmetric hourglass around the latent bottleneck.
    \def\wideedge{1.5}   % half-width of the wide (data) edge
    \def\narrowedge{0.7} % half-width of the narrow (latent) edge
    \def\blockh{1.5}     % block height
    \def\gap{1.2}        % gap between encoder and decoder, sized to fit the latent block

    % Encoder: wide input (bottom) narrows to the code (top)
    \coordinate (e-bl) at (-\wideedge,0);
    \coordinate (e-br) at (\wideedge,0);
    \coordinate (e-tr) at (\narrowedge,\blockh);
    \coordinate (e-tl) at (-\narrowedge,\blockh);
    \draw[thick, fill=rnnblock] (e-bl) -- (e-br) -- (e-tr) -- (e-tl) -- cycle;
    \node[blocklabel] at (0,0.5*\blockh) {Encoder};

    % Decoder: mirrored, the code (bottom) widens to the reconstruction (top)
    \coordinate (d-bl) at (-\narrowedge,\blockh+\gap);
    \coordinate (d-br) at (\narrowedge,\blockh+\gap);
    \coordinate (d-tr) at (\wideedge,2*\blockh+\gap);
    \coordinate (d-tl) at (-\wideedge,2*\blockh+\gap);
    \draw[thick, fill=rnnblock] (d-bl) -- (d-br) -- (d-tr) -- (d-tl) -- cycle;
    \node[blocklabel] at (0,\blockh+\gap+0.5*\blockh) {Decoder};

    % Endpoints of the main data-flow arrow
    \coordinate (xpt) at (0,-1);
    \coordinate (zbot) at (0,\blockh);
    \coordinate (ztop) at (0,\blockh+\gap);
    \coordinate (xhatpt) at (0,2*\blockh+\gap+1);

    % Latent space block, sitting in the bottleneck between the two networks
    \coordinate (zmid) at ($(zbot)!0.5!(ztop)$);
    \node[latentblock] (zblock) at (zmid) {$z$};
    \draw[flow] (xpt) -- (0,0);
    \draw[flow] (zbot) -- (zblock.south);
    \draw[flow] (zblock.north) -- (ztop);
    \draw[flow] (0,2*\blockh+\gap) -- (xhatpt);

    % Endpoint labels
    \node[lbl, below=1mm] at (xpt) {$x$};
    \node[lbl, above=1mm] at (xhatpt) {$\hat{x}$};

    % Side annotations, aligned with the relevant height
    \node[lbl, anchor=east] at ($(xpt)+(-2.4,0)$) {Input data};
    \node[lbl, anchor=east] at ($(zmid)+(-2.4,0)$) {Latent space};
    \node[lbl, anchor=east] at ($(xhatpt)+(-2.4,0)$) {Reconstruction};

    % Reconstruction loss, comparing x and x-hat, routed clear of both blocks
    \draw[skip] (xpt) -- ++(\wideedge+1.0,0) coordinate (sk)
                |- (xhatpt)
                node[skipcolor, midway, right, align=center, font=\small]
                {reconstruction loss\\$\mathcal{L}(x,\hat{x}) = \|x-\hat{x}\|^2$};
\end{tikzpicture}
}
    \caption{Autoencoder architecture: the encoder compresses the input $x$ into a latent code $z$, and the decoder reconstructs $\hat{x}$ from it.}
    \label{fig:autoencoder-architecture}
\end{figure}

To train an autoencoder, we minimize a \textbf{reconstruction loss} between the input and the output of the network.

Formally, given an input $x$, the encoder $f_\theta$ maps it to a latent representation $z = f_\theta(x)$, and the decoder $g_\phi$ reconstructs the input as $\hat{x} = g_\phi(z)$.

A common choice is the Mean Squared Error (MSE):

\begin{equation}
    \mathcal{L}(x,\hat{x}) = \|x-\hat{x}\|^2.
\end{equation}

Depending on the data type, other reconstruction losses may also be used. For example, binary cross-entropy is often preferred when modeling binary or normalized image data.

\paragraph{Implementation}
Autoencoders give us a lot of flexibility in terms of architecture and design choices. In fact, while the idea is fixed, 
the implementation can vary significantly depending on the specific use case and data type.


\paragraph{Why AEs?}
Once trained, the decoder is often discarded and the encoder is used as a feature extractor or for dimensionality reduction. The learned latent representation can then be exploited for a variety of downstream tasks, including visualization, anomaly detection, clustering, and as input features for other machine learning models.

However, standard autoencoders do not explicitly learn a probability distribution over the latent space. As a consequence, although they reconstruct existing samples well, they are not naturally suited for generating new ones.


\subsection{Variational Autoencoders}

Building upon the idea of autoencoders, we obtain a more powerful generative model called the \textbf{Variational Autoencoder} (VAE). Unlike standard autoencoders, VAEs learn not only a latent representation of the data but also a probability distribution over that latent space, making it possible to generate entirely new samples.

VAEs belong to the class of \textbf{explicit density estimation} generative models, whose goal is to directly model the probability distribution of the observed data.

\subsubsection{Idea}
VAEs define a density function $p_\theta(x)$ over the data space, parameterized by a neural network with parameters $\theta$. 

We model the data distribution through

\begin{equation}
    \label{eq:VAEs}
    p_\theta(x)
    =
    \int p(z)\,p_\theta(x|z)\,dz,
\end{equation}

where $p(z)$ is a prior distribution over the latent variable $z$, typically chosen as a standard Gaussian.

Our objective is to maximize the likelihood of the observed data with respect to the parameters $\theta$. Unfortunately, this is generally intractable. Since the conditional distribution $p_\theta(x|z)$ is represented by a neural network, the integral in Equation~\ref{eq:VAEs} has no closed-form solution and must integrate over all possible latent values.

A natural idea is to instead work with the posterior distribution over the latent variables. By Bayes' theorem,

\begin{equation}
    p_\theta(z|x)
    =
    \frac{p_\theta(x|z)\,p(z)}
         {p_\theta(x)}.
\end{equation}

However, this expression still depends on the marginal likelihood $p_\theta(x)$, which is exactly the intractable quantity we wish to compute.

To overcome this difficulty, we introduce a second distribution,

\[
q_\phi(z|x),
\]

parameterized by another neural network. This distribution serves as a tractable approximation of the true posterior $p_\theta(z|x)$.

\paragraph{A tractable solution}

By introducing the approximate posterior $q_\phi(z|x)$ and rearranging the log-likelihood, we obtain the following decomposition:

\begin{equation}
    \log p_\theta(x) = \underbrace{ \mathbb{E}_{q_\phi(z|x)} \left[ \log p_\theta(x|z) \right] - D_{KL} \left(q_\phi(z|x)\,\|\,p(z)\right)}_{\text{ELBO}}+D_{KL}\left(q_\phi(z|x)\,\|\,p_\theta(z|x)\right).
\end{equation}

The three terms have the following interpretation:

\begin{itemize}
    \item The first term is the expected log-likelihood of the data under the approximate posterior. It encourages the decoder to accurately reconstruct the input.

    \item The second term measures how closely the approximate posterior matches the prior distribution over the latent variables.

    \item The third term measures the discrepancy between the approximate and true posterior. Although this quantity is intractable, it is always non-negative.
\end{itemize}



\paragraph{Evidence Lower Bound (ELBO)}

To overcome the intractability of the marginal likelihood, VAEs maximize a tractable lower bound on the log-likelihood, known as the \textbf{Evidence Lower Bound} (ELBO):

\begin{equation}
    \label{eq:ELBO}
    \mathcal{L}_{\mathrm{ELBO}}
    =
    \mathbb{E}_{q_\phi(z|x)}
    \left[
        \log p_\theta(x|z)
    \right]
    -
    D_{KL}
    \left(
        q_\phi(z|x)
        \,\|\,p(z)
    \right).
\end{equation}

The ELBO consists of two terms:

\begin{itemize}
    \item The first term is the \textit{reconstruction term}. It measures how well the decoder $p_\theta(x|z)$ can reconstruct the observed data $x$ from a latent representation sampled from the approximate posterior $q_\phi(z|x)$. Maximizing this term encourages accurate reconstructions.

    \item The second term is the \textit{regularization term}. It is the Kullback--Leibler (KL) divergence between the approximate posterior $q_\phi(z|x)$ and the prior distribution $p(z)$. Minimizing this term encourages the latent representations to remain close to the prior, resulting in a smooth and well-structured latent space.
\end{itemize}

This loss function allows VAEs to learn a meaningful latent space, where similar inputs are mapped to nearby points in the latent space.

\subsubsection{Mapping the ELBO to the VAE Architecture}
Equation \ref{eq:ELBO} maps directly onto the encoder-decoder pair we set out to build, with each network now standing for one of the two distributions involved.

The term $p_\theta(x \mid z)$ represents the probability of generating $x$ given a latent variable $z$. This is exactly what the \textbf{decoder} computes: given a latent code $z$, it outputs a distribution over reconstructions of $x$.

Symmetrically, the term $q_\phi(z \mid x)$ represents the (approximate) probability of a latent variable $z$ given the observed data $x$. This is exactly what the \textbf{encoder} computes: given an input $x$, it outputs a distribution over latent codes $z$, standing in for the true, intractable posterior $p_\theta(z \mid x)$.

Since we are modeling the probabilistic generation of data, both networks are therefore \textbf{probabilistic}: the encoder does not output a single point $z$, but the parameters of the distribution $q_\phi(z \mid x)$ (a mean $\mu_\phi(x)$ and a standard deviation $\sigma_\phi(x)$), from which $z$ is sampled. The decoder, in turn, does not output a single reconstruction, but the parameters of the distribution $p_\theta(x \mid z)$, from which $\hat{x}$ is sampled. This is illustrated in Figure \ref{fig:vae-architecture}.

\begin{figure}[H]
    \centering
    \resizebox{0.5\textwidth}{!}{\usetikzlibrary{arrows.meta, calc}

\begin{tikzpicture}[
    >=Stealth,
    flow/.style={-{Stealth[length=6pt]}, line width=1pt},
    lbl/.style={font=\small, align=center},
    blocklabel/.style={font=\bfseries\small},
    latentblock/.style={rectangle, draw, thick, fill=dlred!20, minimum width=1.4cm, minimum height=0.7cm, font=\small, align=center},
]
    % Geometry, matching chapters/13-autoencoders/tikz/autoencoder-architecture.tex
    % so the two figures read as a single visual family.
    \def\wideedge{1.5}   % half-width of the wide (data) edge
    \def\narrowedge{0.7} % half-width of the narrow (latent) edge
    \def\blockh{1.5}     % block height
    \def\gap{1.2}        % gap between encoder and decoder, sized to fit the latent block

    % Encoder: wide input (bottom) narrows to the approximate posterior q(z|x)
    \coordinate (e-bl) at (-\wideedge,0);
    \coordinate (e-br) at (\wideedge,0);
    \coordinate (e-tr) at (\narrowedge,\blockh);
    \coordinate (e-tl) at (-\narrowedge,\blockh);
    \draw[thick, fill=dlblue!55] (e-bl) -- (e-br) -- (e-tr) -- (e-tl) -- cycle;
    \node[blocklabel] at (0,0.5*\blockh) {$q(z \mid x)$};

    % Decoder: mirrored, the code (bottom) widens to the likelihood p(x|z)
    \coordinate (d-bl) at (-\narrowedge,\blockh+\gap);
    \coordinate (d-br) at (\narrowedge,\blockh+\gap);
    \coordinate (d-tr) at (\wideedge,2*\blockh+\gap);
    \coordinate (d-tl) at (-\wideedge,2*\blockh+\gap);
    \draw[thick, fill=rnnblock] (d-bl) -- (d-br) -- (d-tr) -- (d-tl) -- cycle;
    \node[blocklabel] at (0,\blockh+\gap+0.5*\blockh) {$p(x \mid z)$};

    % Endpoints of the main data-flow arrow
    \coordinate (xpt) at (0,-1);
    \coordinate (zbot) at (0,\blockh);
    \coordinate (ztop) at (0,\blockh+\gap);
    \coordinate (xhatpt) at (0,2*\blockh+\gap+1);

    % Latent space block, sitting in the bottleneck between the two networks
    \coordinate (zmid) at ($(zbot)!0.5!(ztop)$);
    \node[latentblock] (zblock) at (zmid) {$z$};
    \draw[flow] (xpt) -- (0,0);
    \draw[flow] (zbot) -- (zblock.south);
    \draw[flow] (zblock.north) -- (ztop);
    \draw[flow] (0,2*\blockh+\gap) -- (xhatpt);

    % Prior annotation, next to the latent block
    \node[lbl, anchor=west, text=dlred!70!black] at ($(zblock.east)+(0.2,0)$) {$p(z)$};

    % Endpoint labels
    \node[lbl, below=1mm] at (xpt) {$x$};
    \node[lbl, above=1mm] at (xhatpt) {$\hat{x}$};

    % Side annotations, aligned with the relevant height
    \node[lbl, anchor=east] at ($(xpt)+(-2.4,0)$) {Input data};
    \node[lbl, anchor=east] at ($(zmid)+(-2.4,0)$) {Latent space};
    \node[lbl, anchor=east] at ($(xhatpt)+(-2.4,0)$) {Reconstruction};
\end{tikzpicture}
}
    \caption{VAE architecture: the encoder computes the approximate posterior $q_\phi(z \mid x)$, from which a latent code $z$ is sampled, and the decoder computes the likelihood $p_\theta(x \mid z)$, from which the reconstruction $\hat{x}$ is obtained. The prior $p(z)$ regularizes the latent space towards a standard Gaussian.}
    \label{fig:vae-architecture}
\end{figure}

To keep the sampling of $z$ differentiable during training, VAEs use the \textbf{reparameterization trick}: instead of sampling $z$ directly, we sample noise $\epsilon \sim \mathcal{N}(0, I)$ and compute $z = \mu_\phi(x) + \sigma_\phi(x) \odot \epsilon$, moving the randomness outside the network so gradients can still flow through $\mu_\phi$ and $\sigma_\phi$.

\subsubsection{Inference Time}

Unlike a standard autoencoder, a VAE encourages the latent representations to follow a known prior distribution, typically a standard Gaussian. 
As a result, once training is complete, we can generate new data by simply sampling a latent vector

\[
z \sim p(z)
\]

and passing it through the decoder. This regularized latent space is what gives VAEs their generative capability.


\subsubsection{Conditional VAEs}
Similarly to how we can condition GANs on additional information, we can also condition VAEs on auxiliary variables. 
This leads to the \textbf{Conditional Variational Autoencoder} (CVAE), which allows us to generate data conditioned on specific attributes.
This is achieved by concatenating the conditioning variable $c$ to both the input of the encoder and the latent code fed into the decoder, effectively learning a conditional distribution $p_\theta(x|z,c)$.

\subsection{GAN vs VAE}
Even though both GANs and VAEs are generative models, they use very different approaches to generate new data, leading to different strengths and weaknesses.

\begin{itemize}
    \item \textbf{GANs}: provide a mechanism to generate samples, encouraging the model to produce progressively more realistic outputs. 
    However, they do not provide an explicit density function for the data distribution, making it difficult to evaluate the likelihood of generated samples.
    
    \item \textbf{VAEs}: provide an explicit density function for the data distribution. The distribution can then be used both for generating new samples and for evaluating the likelihood of existing samples. 
    However, VAEs often produce blurrier outputs compared to GANs, as they optimize a lower bound on the log-likelihood rather than directly optimizing for sample quality.
\end{itemize}

We can more analitically compare the two models by looking at their foundational properties:

\begin{itemize}
    \item \textbf{Efficient Sampling:} at inference time, both GANs and VAEs can generate new samples efficiently by sampling from their respective latent spaces;
    \item \textbf{Sample Quality:} GANs are known for producing high-quality, realistic samples, due to the adversarial training process that encourages the generator to produce outputs indistinguishable from real data. 
    VAEs, on the other hand, often produce blurrier samples due to the regularization imposed by the KL divergence term in their loss function;
    \item \textbf{Coverage:} VAEs are designed to cover the entire data distribution, as they explicitly model the underlying probability distribution.
    GANs, however, can suffer from mode collapse, where the generator produces a limited variety of samples, failing to capture the full diversity of the data;
    \item \textbf{Well-behaved latent space:} GAN were shown to have a well-structured latent space, allowing for arithmetic operations. 
    For VAEs, the latent space is regularized to follow a known prior distribution, which can also lead to meaningful interpolations between samples.
    \item \textbf{Likelihood Estimation:} VAEs provide a tractable lower bound on the log-likelihood of the data, allowing for efficient likelihood estimation.
    GANs, on the other hand, do not provide a direct way to compute the likelihood of generated samples, making it challenging to evaluate their performance in terms of likelihood.
\end{itemize}